\documentclass[11pt,a4paper]{article}
\usepackage[margin=2.4cm]{geometry}
\usepackage{fontspec}
\setmainfont{Times New Roman}
\usepackage{amsmath,amssymb,empheq}
\usepackage{graphicx,booktabs,caption,enumitem,xcolor}
\usepackage[colorlinks=true,linkcolor=blue!50!black,urlcolor=blue!45!black]{hyperref}
\captionsetup{font=small,labelfont=bf}
\definecolor{boxbg}{HTML}{F4F6F8}\definecolor{warnbg}{HTML}{FFF3CD}\definecolor{keybg}{HTML}{E4EFE4}
\newcommand{\yn}[1]{\par\vspace{1.5mm}\noindent\colorbox{boxbg}{\parbox{0.965\textwidth}{\vspace{1.2mm}#1\vspace{1.2mm}}}\vspace{1.5mm}}
\newcommand{\warn}[1]{\par\vspace{1.5mm}\noindent\colorbox{warnbg}{\parbox{0.965\textwidth}{\vspace{1.2mm}#1\vspace{1.2mm}}}\vspace{1.5mm}}
\newcommand{\key}[1]{\par\vspace{1.5mm}\noindent\colorbox{keybg}{\parbox{0.965\textwidth}{\vspace{1.2mm}#1\vspace{1.2mm}}}\vspace{1.5mm}}

\title{\textbf{Đề cương viết lại}\\[2mm]
\Large Bộ nhớ liên tưởng dày học từ attractor của chính nó:\\
trôi từ đặc trưng sang nguyên mẫu\\[2mm]
\large \normalfont Thiết kế thí nghiệm, đối chứng, và tiêu chí bác bỏ ghi trước}
\author{}\date{}

\begin{document}\maketitle\vspace{-8mm}

\begin{abstract}\noindent
Krotov \& Hopfield (2016) chỉ ra rằng số mũ tách bạch $n$ quyết định một mạng DAM
học ra \emph{đặc trưng} (nét) hay \emph{nguyên mẫu} (chữ số nguyên vẹn) --- một
vị trí \emph{tĩnh} trên một trục, do siêu tham số đặt, với dữ liệu cố định. Đề
cương này đặt câu hỏi mà setting đó không thể hỏi: điều gì xảy ra khi \emph{dữ
liệu chính là đầu ra của mạng}? Giả thuyết trung tâm: bộ ký ức khi đó
\textbf{trôi} dọc đúng trục ấy, luôn về phía nguyên mẫu, và $n$ chỉ quyết định
\emph{tốc độ}, không quyết định \emph{đích}. Hệ quả: năng lực là một \emph{giai
đoạn quá độ}, tăng khi bộ ký ức còn ở vùng đặc trưng và giảm khi nó ngưng tụ.
Chúng tôi đã đo được đúng hình dạng đó trên một hệ spin $150$ biến, nơi biết
được đáp án đúng; đề cương này chuyển nó sang MNIST --- benchmark của chính
Krotov --- với bốn đối chứng và tiêu chí bác bỏ ghi trước khi chạy.
\end{abstract}

%======================================================================
\section{Ghi rõ: cái gì đã có, cái gì là của ta}

\subsection{Đã có (không phải đóng góp)}

\begin{enumerate}[leftmargin=7mm,itemsep=1mm]
\item \textbf{DAM và dung lượng.} Krotov \& Hopfield (2016): năng lượng
$E=-\sum_\mu F(\langle\xi^\mu,v\rangle)$ với $F(x)=\max(x,0)^n$ cho dung lượng
$\sim N^{n-1}$, mũ khi $F=\exp$.
\item \textbf{Chuyển tiếp đặc trưng $\to$ nguyên mẫu.} Cùng bài, trên MNIST:
$n$ thấp thì ký ức học được trông như \emph{nét}; $n$ cao thì trông như
\emph{chữ số nguyên vẹn}. Đây là trục khái niệm mà đề cương này dùng.
\item \textbf{Attention là truy hồi Hopfield hiện đại.} Ramsauer et al.\ (2020):
năng lượng log-sum-exp cho luật cập nhật đúng bằng softmax attention.
\item \textbf{Tự mô hình hoá cải thiện tối ưu hoá.} Watson, Buckley \& Mills
(2011): mạng Hopfield học Hebb trên attractor của chính nó giải bài ràng buộc
tốt hơn.
\end{enumerate}

\warn{\textbf{Phải nói thẳng.} Đánh đổi ``trừu tượng hoá đổi lấy đa dạng theo
$n$'' mà chúng tôi đo được \textbf{chính là} chuyển tiếp đặc trưng--nguyên mẫu
của Krotov, đo trong hoàn cảnh khác. Bản thân nó \emph{không mới}, và không được
trình bày như thể mới.}

\subsection{Cầu nối là một đồng nhất thức}

Lấy self-modeling: các ``ký ức'' $\xi^\mu$ là chính các attractor mà hệ tìm ra.
Với $F(x)=x^2$,
\begin{equation}
\sum_\mu \langle\xi^\mu,s\rangle^2
\;=\; s^{\!\top}\Big(\sum_\mu \xi^\mu\xi^{\mu\top}\Big)s
\;=\; s^{\!\top}\Delta W s ,
\end{equation}
đúng bằng luật Hebb của Watson. Nên
\begin{empheq}[box=\fbox]{equation}
\text{tự mô hình hoá (Watson 2011)} \;=\; \text{DAM tự quy chiếu tại } n=2 .
\end{empheq}
Không phải ẩn dụ. M1 là \emph{một điểm} trên trục của Krotov, và câu hỏi ``ngoài
điểm đó thì sao'' là câu hỏi chưa ai đặt.

\subsection{Giả thuyết trung tâm --- cái mới}

\key{\textbf{H.} Trong DAM có giám sát, dữ liệu cố định, nên $n$ quyết định mạng
\emph{ngồi ở đâu} trên trục đặc trưng--nguyên mẫu, và vị trí đó \emph{tĩnh}. Khi
DAM học từ attractor của chính nó, dữ liệu \emph{là} đầu ra của mô hình, nên hệ
\textbf{di chuyển} dọc trục --- luôn về phía nguyên mẫu --- và $n$ chỉ đặt
\emph{tốc độ}, không đặt \emph{đích}.\\[1.5mm]
\textbf{Hệ quả H'.} Năng lực là giai đoạn quá độ: tăng khi bộ ký ức còn ở vùng
đặc trưng, đạt đỉnh, rồi giảm khi ngưng tụ.}

Điều này \emph{không thể} nghiên cứu trong setting có giám sát, vì ở đó MNIST
không đổi. Vòng phản hồi đơn giản là không tồn tại.

%======================================================================
\section{Bằng chứng sơ bộ đã có}

Trên bài ràng buộc modular $N=150$ ($30$ module $\times\,5$ unit), nơi
\emph{biết được} lời giải tối ưu ($-483{,}3$) và sàn tìm kiếm cá thể ($-419{,}5$):

\begin{table}[h]\centering\small
\begin{tabular}{rrrrr}
\toprule
$n$ & lợi ích học được & trừu tượng hoá (đỉnh) & đa dạng cuối & hạng hiệu dụng\\
\midrule
$2$  & $\mathbf{60}$ & $\mathbf{2{,}23}$ & $\mathbf{1}$ & $290$\\
$6$  & $57$ & $1{,}75$ & $2$ & $259$\\
$10$ & $37$ & $1{,}39$ & $15$ & $94$\\
$20$ & $17$ & $1{,}13$ & $\mathbf{34}$ & $\mathbf{4{,}4}$\\
\bottomrule
\end{tabular}
\caption{Ba cột đối nghịch đơn điệu. Trên toàn bộ $60$ cấu hình $(n,\gamma)$:
$\mathrm{corr}(\text{đa dạng giữ được},\ \text{năng lượng cuối}) = +0{,}809$ ---
giữ được đa dạng thì lời giải tệ đi.}
\end{table}

Và tại $n=2$, quỹ đạo có đúng hình dạng H' dự đoán: trừu tượng hoá lên đỉnh
$2{,}23$ rồi rã về $1{,}03$; đa dạng attractor $46\to1$; $\Delta W$ về
\emph{hạng một} ($\lambda_1/\sum\lambda=0{,}50$,
$\cos(\Delta W,ss^{\!\top})=0{,}90$), tức dung lượng hiệu dụng bằng $1$ trong
khi Hopfield cổ điển đáng lẽ chứa $\approx 21$.

\yn{Đây là bằng chứng \emph{gợi ý}, không phải kết quả: một instance bài toán,
một họ seed, và ``trừu tượng hoá'' đo bằng \emph{một} proxy. Nó biện minh cho
việc làm thí nghiệm, không thay thế thí nghiệm.}

%======================================================================
\section{Thiết kế thí nghiệm trên MNIST}

\subsection{Hệ}

Ma trận ký ức $\Xi\in\mathbb{R}^{K\times 784}$, năng lượng
$E(v) = -\sum_{\mu=1}^{K} F\big(\langle\xi^\mu,v\rangle/\sqrt{N}\big)$,
$F(x)=\max(x,0)^n$. Truy hồi bằng động lực giảm năng lượng.

\textbf{Vòng tự quy chiếu.} Khởi tạo $\Xi$ từ $K$ ảnh MNIST thật --- tức
\emph{bắt đầu ở vùng đặc trưng/dữ liệu}. Mỗi vòng: lấy một truy vấn (ảnh MNIST
nhiễu), thả về attractor $a_t$, rồi \emph{thay} một ký ức bằng $a_t$. Lặp.

\subsection{Đo gì}

\begin{enumerate}[leftmargin=7mm,itemsep=1mm]
\item \textbf{Vị trí trên trục đặc trưng--nguyên mẫu.} Với mỗi ký ức $\xi^\mu$,
lấy profile tương quan $c_i=\mathrm{corr}(\xi^\mu,x_i)$ với toàn bộ ảnh huấn
luyện, rồi tính
\begin{equation}
\mathrm{PR}(\xi^\mu) \;=\; \frac{\big(\sum_i c_i^+\big)^2}{\sum_i (c_i^+)^2},
\qquad c^+=\max(c,0).
\end{equation}
$\mathrm{PR}$ thấp $=$ ký ức giải thích \emph{một} ảnh $=$ nguyên mẫu;
$\mathrm{PR}$ cao $=$ giải thích \emph{nhiều} ảnh $=$ đặc trưng. Đây là lượng
hoá trực tiếp trục của Krotov.
\item \textbf{Hạng hiệu dụng} của $\Xi$ (participation ratio của phổ kỳ dị).
\item \textbf{Đa dạng}: số attractor phân biệt trong cửa sổ trượt.
\item \textbf{Năng lực}: khử nhiễu trên tập giữ lại --- làm hỏng một chữ số chưa
thấy, thả, rồi chấm bằng \emph{một bộ phân loại ngoài cố định}. Chấm luôn trên
nhãn thật, không bao giờ trên năng lượng của mô hình.
\item \textbf{Hình ảnh}: hiện chính các ký ức theo thời gian. Nếu H đúng, ở
$n=2$ ta phải \emph{nhìn thấy} chúng nhoè dần về một vệt trung bình.
\end{enumerate}

\subsection{Đối chứng --- phần đã hỏng ở lần trước}

\warn{Đây là mục quan trọng nhất của đề cương. Lần trước (ABDAM) thất bại đúng ở
đây: một thống kê tổng hợp bị hiểu sai, và cơ chế được tuyên bố không có đối
chứng loại trừ lời giải thích tầm thường. Bốn đối chứng dưới đây phải chạy
\emph{cùng lúc} với thí nghiệm chính, không phải bổ sung sau.}

\begin{description}[leftmargin=6mm,style=nextline]
\item[C1 --- không tự quy chiếu (null).] Mỗi vòng thay ký ức bằng một ảnh MNIST
\emph{thật}, không phải attractor. \emph{Phải không trôi.} Nếu nó cũng trôi thì
hiện tượng không liên quan gì tới tự quy chiếu.
\item[C2 --- đổi siêu tham số đơn thuần.] Bất kỳ hiệu ứng nào quan sát được,
kiểm xem một thay đổi đơn giản của $n$ hoặc thang năng lượng có tái tạo được
không. \emph{Bài học $\beta$}: nếu ``cơ chế'' tương đương một phép đổi thang thì
nó không phải cơ chế.
\item[C3 --- lưu trữ ngẫu nhiên.] Thay ký ức bằng \emph{vector ngẫu nhiên} cùng
chuẩn. Tách ``lưu bất cứ thứ gì'' khỏi ``lưu đầu ra của chính mình''. (Trên hệ
spin, đối chứng tương ứng đã cho kết quả quyết định: nhiễu đối xứng ngẫu nhiên
pha loãng $1{,}00\to0{,}77$ và \emph{không} cải thiện gì, trong khi Hebb thật
$1{,}00\to0{,}11$ và cải thiện $66$ đơn vị.)
\item[C4 --- có neo.] $\Xi \leftarrow (1-\alpha)\Xi_0 + \alpha\Xi + \text{mới}$.
Phải trôi \emph{chậm hơn}, đơn điệu theo $\alpha$.
\end{description}

\subsection{Quy tắc tự vệ, rút từ lần hỏng trước}

\begin{enumerate}[leftmargin=7mm,itemsep=1mm]
\item \textbf{Mọi thống kê tổng hợp phải kiểm hiện vật.} Trước khi diễn giải một
chỉ số bất bình đẳng, tính giá trị của nó cho phân bố \emph{tầm thường} tương
ứng. (Gini của ``$k$ ô bằng nhau $+$ $m$ ô rỗng'' là $1-k/(k+m)$ --- nếu số đo
trùng con số đó thì nó không đo tự tổ chức.)
\item \textbf{Chạy quá đỉnh.} Không dừng ở epoch mà hiệu năng tốt nhất. Chạy
ít nhất $5\times$ thời gian tới đỉnh, để thấy phần sau.
\item \textbf{Nhiều seed và nhiều tập con dữ liệu}, báo cáo sai số chuẩn. Một
instance không phải kết quả.
\item \textbf{Chấm trên đại lượng ngoài.} Không bao giờ dùng năng lượng của
chính mô hình làm thước đo năng lực.
\end{enumerate}

%======================================================================
\section{Tiêu chí bác bỏ --- ghi TRƯỚC khi chạy}

\warn{\begin{itemize}[nosep,leftmargin=6mm]
\item \textbf{H chết} nếu $\mathrm{PR}$ trung bình \emph{không} giảm đơn điệu
theo vòng ở bất kỳ $n$ nào trong $\{2,4,10,20\}$; hoặc nếu nó giảm ở $n$ thấp
nhưng \emph{tăng} ở $n$ cao (khi đó $n$ quyết định đích, không chỉ tốc độ).
\item \textbf{H chết} nếu C3 (lưu ngẫu nhiên) cho độ trôi bằng thí nghiệm chính
trong sai số $20\%$ --- khi đó tự quy chiếu không đóng vai trò gì.
\item \textbf{H chết} nếu C1 (dữ liệu thật) cũng trôi.
\item \textbf{H' chết} nếu năng lực khử nhiễu đơn điệu (chỉ tăng, hoặc chỉ
giảm), không có đỉnh, trên $\ge 3/4$ số $n$.
\item \textbf{Đóng góp bị hạ cấp} nếu mọi hiệu ứng tái tạo được bằng C2.
\end{itemize}
Nếu H chết mà H' sống, vẫn còn một bài (cửa sổ năng lực trong DAM, đo trên
benchmark của Krotov). Nếu cả hai chết thì dừng và không viết gì.}

%======================================================================
\section{Chi phí và rủi ro}

\textbf{Chi phí.} MNIST $784$ chiều, $K\sim100$ ký ức, mỗi vòng là vài phép nhân
ma trận. Hàng nghìn vòng trong vài phút trên CPU. Toàn bộ lưới
$n\times\alpha\times$ seed nằm trong vài giờ. \emph{Không cần GPU.}

\textbf{Rủi ro, xếp theo mức độ nghiêm trọng.}
\begin{enumerate}[leftmargin=7mm,itemsep=1.5mm]
\item \textbf{Hiện tượng có thể tầm thường.} Lưu attractor của chính mình thì
attractor được củng cố --- có thể ai cũng thấy ngay. Phản biện: nếu tầm thường
thì tự mô hình hoá của Watson đã không hoạt động; nội dung định lượng nằm ở
\emph{cửa sổ}, ở biên của nó, và ở việc $n$ điều khiển tốc độ chứ không điều
khiển đích. Nhưng phải trình bày cho đúng tầm.
\item \textbf{Thước đo $\mathrm{PR}$ có thể không sạch.} Cần kiểm trên các
trường hợp đã biết (ảnh thật $\Rightarrow$ PR thấp; ảnh trung bình $\Rightarrow$
PR cao) trước khi tin.
\item \textbf{Có thể đã có người làm.} Tìm kiếm sơ bộ không thấy ``DAM học từ
attractor của chính nó'', nhưng chưa vét hết; phải đọc kỹ mảng động lực học của
DAM (ví dụ arXiv:2506.00851) trước khi tuyên bố tính mới.
\item \textbf{Quy mô.} MNIST với $n$ cao có thể quá dễ, giống như bài toán spin
tự tầm thường hoá khi module lớn. Cần kiểm độ khó trước: nếu khử nhiễu đạt trần
ngay từ vòng đầu thì không có chỗ cho cửa sổ xuất hiện.
\end{enumerate}

%======================================================================
\section{Việc cần làm, theo thứ tự}

\begin{enumerate}[leftmargin=7mm,itemsep=1mm]
\item Kiểm thước đo $\mathrm{PR}$ trên các trường hợp đã biết. Vài phút.
Không qua thì dừng.
\item Kiểm độ khó: khử nhiễu có trần không, có chỗ cho cải thiện không.
\item Thí nghiệm chính $+$ bốn đối chứng, chạy đồng thời, $n\in\{2,4,10,20\}$,
$\ge 5$ seed, chạy quá đỉnh $5\times$.
\item Chỉ khi (3) qua tiêu chí bác bỏ mới viết.
\end{enumerate}

\vspace{3mm}
\noindent\rule{\textwidth}{0.4pt}\\[1mm]
{\small Bằng chứng sơ bộ: \texttt{dam\_selfmodel.py}, \texttt{run\_dam\_sweep.py},
\texttt{data/dam\_sweep.npz}. Hệ spin: $N=150$, $30\times5$, $p=0{,}8$, chấm trên
$A$ với đường chéo bằng không.}

\end{document}
